\documentclass[12pt,a4paper]{book}
%\documentclass[a4paper,french]{book}
%\usepackage[17pt]{extsizes}


\input{../1ere-preambule}

\newcommand{\va}[1]{\lvert \ #1 \ \rvert}
\newcommand{\vaf}[1]{\left\lvert \ #1\  \right\rvert}

\begin{document}
\graphicspath{{images/}}
\setcounter{chapter}{2}
\chapter[ Dérivation - Exercices]{Dérivation \\ Exercices d'entrainement}
\thispagestyle{fancy}

\exerciceDS{}

\begin{minipage}[c]{0.48\linewidth}
	A l'aide du graphique ci-contre et sachant que la courbe verte est la courbe représentative de la fonction f et que les droites rouges sont les tangentes à la courbe en 0,2 et 6, déterminer $f'(2)$, $f'(0)$ et $f'(6)$.
\end{minipage} \hfill
\begin{minipage}[c]{0.48\linewidth}
	\begin{center}
		\newrgbcolor{qqwuqq}{0. 0.39215686274509803 0.}
		\psset{xunit=0.8cm,yunit=0.8cm,algebraic=true,dimen=middle,dotstyle=o,dotsize=5pt 0,linewidth=1.5pt,arrowsize=3pt 2,arrowinset=0.25}
		\begin{pspicture*}(-2.5,-5.)(7.,5.2)
			\multips(0,-5)(0,1.0){11}{\psline[linestyle=dashed,linecap=1,dash=1.5pt 1.5pt,linewidth=0.4pt,linecolor=lightgray]{c-c}(-2.5,0)(7.,0)}
			\multips(-2,0)(1.0,0){10}{\psline[linestyle=dashed,linecap=1,dash=1.5pt 1.5pt,linewidth=0.4pt,linecolor=lightgray]{c-c}(0,-5.)(0,5.2)}
			\psaxes[labelFontSize=\scriptstyle,xAxis=true,yAxis=true,Dx=1.,Dy=1.,ticksize=-2pt 0,subticks=2]{->}(0,0)(-2.5,-5.)(7.,5.2)
			\psplot[linewidth=1.2pt,linecolor=qqwuqq,plotpoints=200]{-2.5}{7.0}{-0.5*x^(2.0)+2.0*x+2.0}
			\psplot[linewidth=1.2pt,linecolor=red]{-2.5}{7.}{(--2.--2.*x)/1.}
			\psplot[linewidth=1.2pt,linecolor=red]{-2.5}{7.}{(--4.-0.*x)/1.}
			\psplot[linewidth=1.2pt,linecolor=red]{-2.5}{7.}{(--20.-4.*x)/1.}
			%		\begin{scriptsize}
				\rput[bl](-3.52,-11.04){\qqwuqq{$f$}}
				\rput[bl](2.22,5.82){\red{$g$}}
				\rput[bl](-4.14,3.68){\red{$h$}}
				\rput[bl](3.62,5.82){\red{$i$}}
				%		\end{scriptsize}
		\end{pspicture*}
	\end{center}
\end{minipage}





\exerciceDS{}

\begin{minipage}[c]{0.48\linewidth}
	Dans le graphique ci-contre, tracer les tangentes en 1 et en 3 à la courbe représentative de la fonction f sachant que $f'(1)=-1$ et $f'(3)=3$
\end{minipage} \hfill
\begin{minipage}[c]{0.48\linewidth}
	\begin{center}
		\psset{xunit=0.8cm,yunit=0.8cm,algebraic=true,dimen=middle,dotstyle=o,dotsize=5pt 0,linewidth=1.2pt,arrowsize=3pt 2,arrowinset=0.25}
		\begin{pspicture*}(-2.,-3.5)(5.,4.5)
			\multips(0,-3)(0,1.0){9}{\psline[linestyle=dashed,linecap=1,dash=1.5pt 1.5pt,linewidth=0.4pt,linecolor=lightgray]{c-c}(-2.,0)(5.,0)}
			\multips(-2,0)(1.0,0){8}{\psline[linestyle=dashed,linecap=1,dash=1.5pt 1.5pt,linewidth=0.4pt,linecolor=lightgray]{c-c}(0,-3.5)(0,4.5)}
			\psaxes[labelFontSize=\scriptstyle,xAxis=true,yAxis=true,Dx=1.,Dy=1.,ticksize=-2pt 0,subticks=2]{->}(0,0)(-2.,-3.5)(5.,4.5)
			\psplot[linewidth=1.2pt,linecolor=ForestGreen,plotpoints=200]{-2.0}{5.0}{x^(2.0)-3.0*x}
		\end{pspicture*}
	\end{center}
\end{minipage}

\newpage

\exerciceDS{}

Déterminer sur $\mathbb{R}$ l'expression de la fonction dérivée de chacune des fonctions suivantes.
\begin{enumMet}
	\item $f(x)= 2x^2-5 x-5$
	\item $g(x)= 6 x^3-3 x^2+7$
	\item $h(x)= 2x^3-5x^2-4x+1$
	\item $k(x)=-7 x^3-6 x-4x^2+3$
\end{enumMet}

\exerciceDS{}

Déterminer sur $\mathbb{R}$ l'expression de la fonction dérivée de chacune des fonctions suivantes.
\begin{enumMet}
	\item $f(x)= 0,7 x^2+6,3 x-5$
	\item $g(x)= 2-8 x^2$
	\item $h(x)= 4 x^3+16 x-15$
	\item $k(x)=-\dfrac{7}{3} x^3+\dfrac{9}{2} x^2$
\end{enumMet}

\exerciceDS{}

Soit $f$ la fonction définie sur \R par $f(x)=-2x^2+8x-5$
\begin{enumMet}
	\item Déterminer $f'$ la fonction dérivée de $f$ sur \R
	\item Etudier le signe de $f'(x)$ sur \R et en déduire les variations de la fonction $f$ sur \R
\end{enumMet}

\exerciceDS{}

Soit $f$ la fonction définie sur \R par $f(x)=-2x^3+\dfrac{25}{2}x^2-4x-1$
\begin{enumMet}
	\item Déterminer $f'$ la fonction dérivée de $f$ sur \R
	\item Montrer que $f'(x)=(-6x+1)(x-4)$
	\item Etudier le signe de $f'(x)$ sur \R et en déduire les variations de la fonction $f$ sur \R
\end{enumMet}





\end{document}

